%% Postdoctoral Research Statement
%% Adapted from the SOP template at
%% http://tex.stackexchange.com/questions/150900/latex-coding-for-statement-of-purpose
%%
%% Compile with:  pdflatex -> bibtex -> pdflatex -> pdflatex
%% Requires references.bib alongside this file.

\documentclass[12pt]{article}
\renewcommand{\baselinestretch}{1} 
\usepackage[
  a4paper,
  margin=1in,
  headsep=3pt, % separation between header rule and text
]{geometry}
\usepackage[most]{tcolorbox}
\setlength{\floatsep}{3pt plus 1pt minus 1pt}
\setlength{\textfloatsep}{3pt plus 1pt minus 1pt}
\setlength{\intextsep}{3pt plus 1pt minus 1pt}
\usepackage{xcolor}
\usepackage{fancyhdr}
\usepackage{scholax}
\usepackage{lastpage}
\usepackage{lineno}
% \linenumbers
\usepackage[inline]{enumitem}
\usepackage{caption}
\usepackage{wrapfig}
\usepackage{graphicx}
\usepackage{ragged2e}
\usepackage{blindtext}
\usepackage{subcaption}
\usepackage{xcolor}
\usepackage{marginnote}
\usepackage{caption}
\usepackage{enumitem}
\setlist[itemize]{noitemsep}
\newcommand{\tcolor}{red}
\DeclareCaptionFont{customsizefig}{\fontsize{10pt}{10pt}\selectfont}
\DeclareCaptionFont{customsizesubfig}{\fontsize{9pt}{9pt}\selectfont}
\captionsetup[figure]{font=customsizefig,aboveskip=1.5pt, belowskip=1.5pt}
\captionsetup[subfigure]{font=customsizesubfig,aboveskip=1.5pt, belowskip=1.5pt}
\newcommand{\TODO}[1]{%
  \marginnote{\color{red}\textbf{TODO}}% Put 'TODO' in margin
  {\color{red}#1}% Put the associated text in red in the body
}
\setlist[itemize]{itemsep=-0.1cm, topsep=0pt, leftmargin=1.2em}
%% ---------------------------------------------------------------
%% Bibliography: natbib + BibTeX, author-year style.  Cite in-line
%% with \citep{key} (parenthetical) or \citet{key} (textual).
%% ---------------------------------------------------------------
\usepackage[round,authoryear]{natbib}
\bibliographystyle{ieeetr}
\setlength{\bibsep}{2pt plus 0.3ex}

\pagestyle{fancy}
\fancyhf{}
\fancyhead[C]{%
  \footnotesize\rmfamily
  Web: \textcolor{blue}{\yourweb}\quad
  \yourname\ \quad
  E-mail: \textcolor{blue}{\youremail}}
\fancyfoot[C]{Page \thepage\ of \pageref{LastPage}}

%% ---------------------------------------------------------------
%% Personal-info macros: edit these once, they propagate everywhere.
%% ---------------------------------------------------------------
\newcommand{\yourname}{Sean Patrick MacBride}
\newcommand{\youremail}{\href{mailto:sean.macbride@physik.uzh.ch}{sean.macbride@physik.uzh.ch}}
\newcommand{\yourweb}{\href{https://seanmacb.github.io}{seanmacb.github.io}}

\newcommand{\hindex}{8}
\newcommand{\hindexDate}{August 2026}

\newcommand{\soptitle}{Publication Impact Statement}

%% Section-header macro (kept from the original SOP template)
\newcommand{\statement}[1]{\par\medskip
  \textcolor{blue}{\textbf{#1:}}%\hspace{0.05cm}
}

%% Highlight macro for spots the author should tailor per application
%% (e.g. naming the specific host group, fellowship, or institution).
\newcommand{\fillin}[1]{\textcolor{orange}{\textbf{[#1]}}}

%% ---------------------------------------------------------------
%% Recurring survey / collaboration / instrument acronyms.
%%
%% These are now managed by the `acro' package: each is declared once
%% with a long and a short form, and \ac{...} automatically prints
%% "Long Form (SHORT)" on first use and "SHORT" every time after.
%% The user-facing macros (\LSST, \DESC, ...) are kept unchanged so
%% the body text below did not have to be rewritten.
%% ---------------------------------------------------------------
\usepackage{acro}
\acsetup{make-links=false}

\DeclareAcronym{lsst}    {short = LSST,    long = Legacy Survey of Space and Time}
\DeclareAcronym{lsstcam} {short = LSSTCam, long = \textbf{LSST Camera}}
\DeclareAcronym{gw} {short = GW, long = gravitational wave}
\DeclareAcronym{em} {short = EM, long = electromagnetic}
\DeclareAcronym{desc}    {short = DESC,    long = LSST-Dark Energy Science Collaboration}
\DeclareAcronym{des}     {short = DES,     long = Dark Energy Survey}
\DeclareAcronym{desi}    {short = DESI,    long = \textbf{Dark Energy Spectroscopic Instrument}}
\DeclareAcronym{desgw}   {short = DESGW,   long = Dark Energy Survey Gravitational Wave follow-up program}
\DeclareAcronym{delve}   {short = DELVE,   long = Dark Energy Local Volume Exploration}
\DeclareAcronym{pfs}     {short = PFS,     long = \textbf{Prime Focus Spectrograph}}
\DeclareAcronym{llamas}  {short = LLAMAS,  long = \textbf{Large Lenslet Array Magellan Spectrograph}}
\DeclareAcronym{lvk}     {short = LVK,     long = LIGO-Virgo-Kagra}
\DeclareAcronym{too}     {short = ToO,     long = Target-of-Opportunity}
\DeclareAcronym{decam}   {short = DECam,   long = \textbf{Dark Energy Camera}}
\DeclareAcronym{rubin}   {short = Rubin,   long = Vera C. Rubin Observatory}
\DeclareAcronym{must}   {short = MUST,   long = MUltiplexed Survey Telescope}
\DeclareAcronym{3G}   {short = 3G,   long = third generation}
\DeclareAcronym{specsfive}   {short = Spec-S5,   long = \textbf{Stage-5 Spectroscopic Experiment}}
\DeclareAcronym{LISA}   {short = LISA,   long = Laser Interferometer Space Antenna}
\DeclareAcronym{hyperk}   {short = Hyper-K,   long = Hyper Kamiokande}
\DeclareAcronym{DUNE}   {short = DUNE,   long = Deep Underground Neutrino Experiment}
\DeclareAcronym{jgem}    {short = J-GEM,
                          long  = Japanese collaboration for gravitational-wave electromagnetic follow-up}
\DeclareAcronym{UNIONS}   {short = UNIONS,   long = The Ultraviolet Near-infrared Optical Northern Survey} 

\newcommand{\LSST}{\ac{lsst}}
\newcommand{\EM}{\ac{em}}
\newcommand{\threeG}{\ac{3G}}
\newcommand{\UNIONS}{\ac{UNIONS}}
\newcommand{\LSSTCam}{\ac{lsstcam}}
\newcommand{\DESC}{\ac{desc}}
\newcommand{\DES}{\ac{des}}
\newcommand{\DESI}{\ac{desi}}
\newcommand{\DESGW}{\ac{desgw}}
\newcommand{\DELVE}{\ac{delve}}
\newcommand{\PFS}{\ac{pfs}}
\newcommand{\LLAMAS}{\ac{llamas}}
\newcommand{\LVK}{\ac{lvk}}
\newcommand{\ToO}{\ac{too}}
\newcommand{\DECam}{\ac{decam}}
\newcommand{\JGEM}{\ac{jgem}}
\newcommand{\GW}{\ac{gw}}
\newcommand{\must}{\ac{must}}

\newcommand{\specsfive}{\ac{specsfive}}
\newcommand{\LISA}{\ac{LISA}}
\newcommand{\hyperk}{\ac{hyperk}}
\newcommand{\DUNE}{\ac{DUNE}}

%% Non-acronym text macros (unchanged: these are names, not abbreviations)
\newcommand{\Rubin}{\ac{rubin}}
\newcommand{\Ho}{\ensuremath{H_0}}
\newcommand{\Blanco}{V\'ictor M. Blanco 4-meter Telescope}
\newcommand{\IceCube}{IceCube Neutrino Observatory}

\newcommand{\Institute}{\textcolor{\tcolor}{Fermilab National Accelerator Laboratory}}
\newcommand{\Position}{\textcolor{\tcolor}{Lederman Fellow}}

\usepackage[
  breaklinks,
  pdftitle={\yourname - \soptitle},
  pdfauthor={\yourname},
  unicode
]{hyperref}

\begin{document}
\hrule
\begin{center}
    \Large\soptitle\ -- \yourname
\end{center}
\vspace{0pt}
\hrule
\vspace{0.5pt}
\hrule height 0.5pt
\vspace{3pt}

\normalsize
\noindent
\thispagestyle{plain}
Three publications demonstrate my technical and analytic expertise from my PhD, spanning astronomical instrumentation development, \ToO\ searches, and standard siren cosmology.

\textbf{The On-Sky Performance of the LSST Camera CCD Array.} 

\textit{Sean P. MacBride, Aaron Roodman, Stuart Marshall, et al. X-‌Ray, Optical, and Infrared Detectors for Astronomy XII 14157, 272-‌292, SPIE.}

In this paper, we report the on-sky performance of LSST-Camera, from arrival at Rubin in May 2024 through its first year of operations in mid-2025. The paper covers recovering CCDs affected by an electrical short and a faulty ADC, retuning readout timing for survey strategy changes, and adding an active clearing safety feature. It also describes on-sky sensor features and overall performance. My contribution was through commissioning and characterizing the LSST-Camera during my time in Chile. The publication impact is from the first on-sky characterization of the largest focal plane ever built.

\textbf{A Joint Search for the Electromagnetic Counterpart to the Gravitational-Wave Binary Black-Hole Merger Candidate S250328ae with the Dark Energy Camera and the Prime Focus Spectrograph.} 

\textit{Haibin Zhang, Mitsuru Kokubo, Sean MacBride, et al. Submitted to The Astrophysical Journal, \href{https://arxiv.org/abs/2508.00291}{\texttt{arXiv:2508.00291}}.} 

In this study, we searched for an optical counterpart to the binary black hole merger candidate S250328ae ($511 \pm 82\ \text{Mpc}$, $90\%$ localization of $15\ \deg^2$). We imaged the localization with DECam, found 36 high-confidence transients, and followed up with PFS spectroscopy of 3,897 targets covering about half the area. No confident counterpart was found, though a few faint candidates remain. I led the DECam observations and transient analysis. This analysis is the first science result from PFS, and a model for combining wide-field imagers with multi-object spectrographs to search for electromagnetic counterparts to gravitational waves.

\textbf{Constraints on the Cosmic Expansion Rate Using the DECam Local Volume Exploration Survey and the Fifth Gravitational Wave Transient Catalog.} 

\textit{Sean Patrick MacBride, Isaac McMahon, Danny Laghi, et al. Submitted to Astronomy \& Astrophysics, \href{https://arxiv.org/abs/2610.25909}{arXiv:2610.25909}.}

In this study, we measure $H_0$ with the dark siren method, combining the DELVE DR3 galaxy catalog with GWTC-5.0. The wider, shallower, photometrically noisier footprint required new catalog choices: a faint magnitude cut ($m_r < 20$), a bright cut ($m_r > 15.5$), and an evolving luminosity function. The result is $H_0 = 67.4^{+13.0}_{-11.2}\ \text{km s}^{-1}\ \text{Mpc}^{-1}$ from dark sirens alone and $H_0 = 70.6^{+8.7}_{-7.0}\ \text{km s}^{-1}\ \text{Mpc}^{-1}$ with GW170817. I led this analysis, including catalog preparation, systematic assessment, and interpretation. It shows the method works on wider-area, non-uniform surveys, and serves as a template for LSST, Euclid and Roman catalogs.

\end{document}
